\documentclass[aspectratio=169, 11pt]{beamer}
% ============================================================================
% Preambulo compartido para presentaciones Beamer
% Estadistica 2026-II - Pontificia Universidad Javeriana
% Incluir con: % ============================================================================
% Preambulo compartido para presentaciones Beamer
% Estadistica 2026-II - Pontificia Universidad Javeriana
% Incluir con: % ============================================================================
% Preambulo compartido para presentaciones Beamer
% Estadistica 2026-II - Pontificia Universidad Javeriana
% Incluir con: \input{../preambulo_beamer.tex} despues de \documentclass
% ============================================================================

\usepackage[T1]{fontenc}
\usepackage[utf8]{inputenc}
\usepackage[spanish]{babel}
\usepackage{amsmath, amssymb, amsfonts}
\usepackage{booktabs}
\usepackage{graphicx}
\usepackage{tikz}
\usetikzlibrary{shapes.geometric, positioning, arrows.meta, calc}
\usepackage{pgfplots}
\pgfplotsset{compat=1.18}
\usepackage{listings}
\usepackage{xcolor}
\usepackage{hyperref}
\usepackage{multicol}

% --- Tema Beamer (minimalista) ---
\usetheme{default}
\usecolortheme{default}
\usefonttheme{default}
\setbeamertemplate{navigation symbols}{}
\setbeamertemplate{caption}[numbered]
\setbeamertemplate{footline}{%
  \hspace{1em}{\tiny\url{https://github.com/polanco-jaime/Curso-Estadistica-2026-3}}\hfill\usebeamerfont{footline}\insertframenumber/\inserttotalframenumber\hspace{1em}\vspace{0.5em}%
}
\setbeamertemplate{itemize items}[circle]
\setbeamertemplate{enumerate items}[default]

% --- Colores institucionales (sutiles) ---
\definecolor{javeriana}{RGB}{0, 62, 105}
\definecolor{javerianalight}{RGB}{0, 105, 170}
\definecolor{codigobg}{RGB}{245, 245, 245}
\definecolor{codigofg}{RGB}{40, 40, 40}
\definecolor{comentario}{RGB}{100, 100, 100}
\definecolor{keyword}{RGB}{0, 100, 150}
\definecolor{string}{RGB}{180, 60, 30}

\setbeamercolor{title}{fg=javeriana}
\setbeamercolor{frametitle}{fg=javeriana}
\setbeamercolor{structure}{fg=javeriana}
\setbeamercolor{block title}{fg=white, bg=javeriana!75}
\setbeamercolor{block body}{bg=javeriana!5}
\setbeamercolor{block title alerted}{fg=white, bg=red!70!black}
\setbeamercolor{block body alerted}{bg=red!5}
\setbeamercolor{block title example}{fg=white, bg=green!40!black}
\setbeamercolor{block body example}{bg=green!5}
\setbeamercolor{item}{fg=javeriana}

\setbeamerfont{title}{series=\bfseries, size=\Large}
\setbeamerfont{frametitle}{series=\bfseries}

% --- Configuracion de listings para R ---
\lstset{
  language=R,
  basicstyle=\ttfamily\footnotesize,
  backgroundcolor=\color{codigobg},
  commentstyle=\color{comentario}\itshape,
  keywordstyle=\color{keyword}\bfseries,
  stringstyle=\color{string},
  numbers=none,
  frame=single,
  rulecolor=\color{javeriana!30},
  breaklines=true,
  showstringspaces=false,
  columns=flexible,
  morekeywords={library, ggplot, aes, geom_histogram, geom_boxplot,
    geom_point, geom_smooth, geom_density, geom_bar, geom_line,
    geom_vline, geom_hline, geom_errorbar, labs, theme_minimal,
    facet_wrap, scale_fill_manual, scale_color_manual,
    summarise, group_by, filter, mutate, select, arrange,
    read.csv, head, str, summary, mean, median, sd, var, cor,
    lm, t.test, chisq.test, wilcox.test, kruskal.test, aov,
    pnorm, qnorm, dnorm, rnorm, qt, pt, qchisq, pchisq, qf, pf,
    confint, coeftest, vcovHC, bptest, vif, cooks.distance,
    TukeyHSD, table, prop.table}
}

% --- Comandos personalizados ---
\newcommand{\R}{\texttt{R}}
\newcommand{\code}[1]{\texttt{\footnotesize #1}}
\newcommand{\variable}[1]{\texttt{#1}}
\newcommand{\paquete}[1]{\texttt{#1}}

% --- Informacion del curso ---
\institute{Pontificia Universidad Javeriana\\
  Facultad de Ciencias Econ\'omicas y Administrativas}
\date{2026-II}
 despues de \documentclass
% ============================================================================

\usepackage[T1]{fontenc}
\usepackage[utf8]{inputenc}
\usepackage[spanish]{babel}
\usepackage{amsmath, amssymb, amsfonts}
\usepackage{booktabs}
\usepackage{graphicx}
\usepackage{tikz}
\usetikzlibrary{shapes.geometric, positioning, arrows.meta, calc}
\usepackage{pgfplots}
\pgfplotsset{compat=1.18}
\usepackage{listings}
\usepackage{xcolor}
\usepackage{hyperref}
\usepackage{multicol}

% --- Tema Beamer (minimalista) ---
\usetheme{default}
\usecolortheme{default}
\usefonttheme{default}
\setbeamertemplate{navigation symbols}{}
\setbeamertemplate{caption}[numbered]
\setbeamertemplate{footline}{%
  \hspace{1em}{\tiny\url{https://github.com/polanco-jaime/Curso-Estadistica-2026-3}}\hfill\usebeamerfont{footline}\insertframenumber/\inserttotalframenumber\hspace{1em}\vspace{0.5em}%
}
\setbeamertemplate{itemize items}[circle]
\setbeamertemplate{enumerate items}[default]

% --- Colores institucionales (sutiles) ---
\definecolor{javeriana}{RGB}{0, 62, 105}
\definecolor{javerianalight}{RGB}{0, 105, 170}
\definecolor{codigobg}{RGB}{245, 245, 245}
\definecolor{codigofg}{RGB}{40, 40, 40}
\definecolor{comentario}{RGB}{100, 100, 100}
\definecolor{keyword}{RGB}{0, 100, 150}
\definecolor{string}{RGB}{180, 60, 30}

\setbeamercolor{title}{fg=javeriana}
\setbeamercolor{frametitle}{fg=javeriana}
\setbeamercolor{structure}{fg=javeriana}
\setbeamercolor{block title}{fg=white, bg=javeriana!75}
\setbeamercolor{block body}{bg=javeriana!5}
\setbeamercolor{block title alerted}{fg=white, bg=red!70!black}
\setbeamercolor{block body alerted}{bg=red!5}
\setbeamercolor{block title example}{fg=white, bg=green!40!black}
\setbeamercolor{block body example}{bg=green!5}
\setbeamercolor{item}{fg=javeriana}

\setbeamerfont{title}{series=\bfseries, size=\Large}
\setbeamerfont{frametitle}{series=\bfseries}

% --- Configuracion de listings para R ---
\lstset{
  language=R,
  basicstyle=\ttfamily\footnotesize,
  backgroundcolor=\color{codigobg},
  commentstyle=\color{comentario}\itshape,
  keywordstyle=\color{keyword}\bfseries,
  stringstyle=\color{string},
  numbers=none,
  frame=single,
  rulecolor=\color{javeriana!30},
  breaklines=true,
  showstringspaces=false,
  columns=flexible,
  morekeywords={library, ggplot, aes, geom_histogram, geom_boxplot,
    geom_point, geom_smooth, geom_density, geom_bar, geom_line,
    geom_vline, geom_hline, geom_errorbar, labs, theme_minimal,
    facet_wrap, scale_fill_manual, scale_color_manual,
    summarise, group_by, filter, mutate, select, arrange,
    read.csv, head, str, summary, mean, median, sd, var, cor,
    lm, t.test, chisq.test, wilcox.test, kruskal.test, aov,
    pnorm, qnorm, dnorm, rnorm, qt, pt, qchisq, pchisq, qf, pf,
    confint, coeftest, vcovHC, bptest, vif, cooks.distance,
    TukeyHSD, table, prop.table}
}

% --- Comandos personalizados ---
\newcommand{\R}{\texttt{R}}
\newcommand{\code}[1]{\texttt{\footnotesize #1}}
\newcommand{\variable}[1]{\texttt{#1}}
\newcommand{\paquete}[1]{\texttt{#1}}

% --- Informacion del curso ---
\institute{Pontificia Universidad Javeriana\\
  Facultad de Ciencias Econ\'omicas y Administrativas}
\date{2026-II}
 despues de \documentclass
% ============================================================================

\usepackage[T1]{fontenc}
\usepackage[utf8]{inputenc}
\usepackage[spanish]{babel}
\usepackage{amsmath, amssymb, amsfonts}
\usepackage{booktabs}
\usepackage{graphicx}
\usepackage{tikz}
\usetikzlibrary{shapes.geometric, positioning, arrows.meta, calc}
\usepackage{pgfplots}
\pgfplotsset{compat=1.18}
\usepackage{listings}
\usepackage{xcolor}
\usepackage{hyperref}
\usepackage{multicol}

% --- Tema Beamer (minimalista) ---
\usetheme{default}
\usecolortheme{default}
\usefonttheme{default}
\setbeamertemplate{navigation symbols}{}
\setbeamertemplate{caption}[numbered]
\setbeamertemplate{footline}{%
  \hspace{1em}{\tiny\url{https://github.com/polanco-jaime/Curso-Estadistica-2026-3}}\hfill\usebeamerfont{footline}\insertframenumber/\inserttotalframenumber\hspace{1em}\vspace{0.5em}%
}
\setbeamertemplate{itemize items}[circle]
\setbeamertemplate{enumerate items}[default]

% --- Colores institucionales (sutiles) ---
\definecolor{javeriana}{RGB}{0, 62, 105}
\definecolor{javerianalight}{RGB}{0, 105, 170}
\definecolor{codigobg}{RGB}{245, 245, 245}
\definecolor{codigofg}{RGB}{40, 40, 40}
\definecolor{comentario}{RGB}{100, 100, 100}
\definecolor{keyword}{RGB}{0, 100, 150}
\definecolor{string}{RGB}{180, 60, 30}

\setbeamercolor{title}{fg=javeriana}
\setbeamercolor{frametitle}{fg=javeriana}
\setbeamercolor{structure}{fg=javeriana}
\setbeamercolor{block title}{fg=white, bg=javeriana!75}
\setbeamercolor{block body}{bg=javeriana!5}
\setbeamercolor{block title alerted}{fg=white, bg=red!70!black}
\setbeamercolor{block body alerted}{bg=red!5}
\setbeamercolor{block title example}{fg=white, bg=green!40!black}
\setbeamercolor{block body example}{bg=green!5}
\setbeamercolor{item}{fg=javeriana}

\setbeamerfont{title}{series=\bfseries, size=\Large}
\setbeamerfont{frametitle}{series=\bfseries}

% --- Configuracion de listings para R ---
\lstset{
  language=R,
  basicstyle=\ttfamily\footnotesize,
  backgroundcolor=\color{codigobg},
  commentstyle=\color{comentario}\itshape,
  keywordstyle=\color{keyword}\bfseries,
  stringstyle=\color{string},
  numbers=none,
  frame=single,
  rulecolor=\color{javeriana!30},
  breaklines=true,
  showstringspaces=false,
  columns=flexible,
  morekeywords={library, ggplot, aes, geom_histogram, geom_boxplot,
    geom_point, geom_smooth, geom_density, geom_bar, geom_line,
    geom_vline, geom_hline, geom_errorbar, labs, theme_minimal,
    facet_wrap, scale_fill_manual, scale_color_manual,
    summarise, group_by, filter, mutate, select, arrange,
    read.csv, head, str, summary, mean, median, sd, var, cor,
    lm, t.test, chisq.test, wilcox.test, kruskal.test, aov,
    pnorm, qnorm, dnorm, rnorm, qt, pt, qchisq, pchisq, qf, pf,
    confint, coeftest, vcovHC, bptest, vif, cooks.distance,
    TukeyHSD, table, prop.table}
}

% --- Comandos personalizados ---
\newcommand{\R}{\texttt{R}}
\newcommand{\code}[1]{\texttt{\footnotesize #1}}
\newcommand{\variable}[1]{\texttt{#1}}
\newcommand{\paquete}[1]{\texttt{#1}}

% --- Informacion del curso ---
\institute{Pontificia Universidad Javeriana\\
  Facultad de Ciencias Econ\'omicas y Administrativas}
\date{2026-II}

\title{Sesión 5: IC para Proporciones y Tamaño de Muestra}
\subtitle{Completando la caja de herramientas de estimación}
\author{Prof.\ Jaime Polanco Jim\'enez}
\date{Viernes 18 de septiembre de 2026}
\begin{document}

\begin{frame}
\titlepage
\end{frame}

\begin{frame}{Agenda}
\tableofcontents
\end{frame}

% ============================================================================
\section{IC para Proporciones}
% ============================================================================

\begin{frame}{Motivación: Estimar Proporciones}
\textbf{Ejemplos de proporciones en educación}:
\begin{itemize}
  \item ¿Qué proporción de estudiantes de NI alcanza nivel B1 en inglés?
  \item ¿Qué porcentaje de estudiantes aprueba el examen ICFES con más de 300 puntos globales?
  \item ¿Cuál es la tasa de deserción en el primer año?
\end{itemize}

\vspace{0.3cm}
\textbf{Diferencia con medias}:
\begin{itemize}
  \item Variable de interés: \alert{categórica binaria} (éxito/fracaso, sí/no)
  \item Codificación: 1 = éxito, 0 = fracaso
  \item Proporción muestral: $\hat{p} = \frac{\text{número de éxitos}}{n} = \frac{x}{n}$
\end{itemize}

\vspace{0.3cm}
\textbf{Objetivo}: Construir IC para la proporción poblacional $p$.
\end{frame}

\begin{frame}{Proporción Muestral $\hat{p}$}
\begin{definition}[Proporción muestral]
\[
\hat{p} = \frac{x}{n}
\]
donde:
\begin{itemize}
  \item $x$ = número de éxitos en la muestra
  \item $n$ = tamaño de la muestra
  \item $\hat{p}$ estima la proporción poblacional $p$
\end{itemize}
\end{definition}

\vspace{0.3cm}
\textbf{Ejemplo}: En una muestra de $n = 200$ estudiantes de NI, 140 superan B1 en inglés.
\[
\hat{p} = \frac{140}{200} = 0.70 = 70\%
\]

\vspace{0.3cm}
\textbf{Conexión con Bernoulli}: Cada estudiante es un ensayo Bernoulli con probabilidad $p$ de éxito. Entonces $x \sim \text{Binomial}(n, p)$.
\end{frame}

\begin{frame}{Distribución Muestral de $\hat{p}$}
\textbf{Propiedades}:
\begin{itemize}
  \item Valor esperado: $E(\hat{p}) = p$ (insesgado)
  \item Varianza: $\text{Var}(\hat{p}) = \frac{p(1-p)}{n}$
  \item Desviación estándar (error estándar): $\text{DE}(\hat{p}) = \sqrt{\frac{p(1-p)}{n}}$
\end{itemize}

\vspace{0.3cm}
\begin{block}{Teorema Central del Límite para proporciones}
Si $n$ es suficientemente grande, entonces:
\[
\hat{p} \sim N\left(p, \frac{p(1-p)}{n}\right) \quad \text{aproximadamente}
\]
Es decir:
\[
Z = \frac{\hat{p} - p}{\sqrt{\frac{p(1-p)}{n}}} \sim N(0,1) \quad \text{aproximadamente}
\]
\end{block}
\end{frame}

\begin{frame}{Condición de Aplicabilidad}
\textbf{¿Cuándo es válida la aproximación normal?}

\vspace{0.3cm}
\begin{alertblock}{Regla empírica}
La aproximación normal es adecuada si:
\[
n \cdot \hat{p} \geq 5 \quad \text{y} \quad n \cdot (1 - \hat{p}) \geq 5
\]
\end{alertblock}

\vspace{0.3cm}
\textbf{Interpretación}:
\begin{itemize}
  \item Necesitamos al menos 5 éxitos y 5 fracasos en la muestra
  \item Garantiza que la distribución binomial no esté muy sesgada
\end{itemize}

\vspace{0.3cm}
\textbf{Ejemplo}: $n = 200$, $\hat{p} = 0.70$
\begin{itemize}
  \item $n \cdot \hat{p} = 200 \times 0.70 = 140 \geq 5$ \checkmark
  \item $n \cdot (1 - \hat{p}) = 200 \times 0.30 = 60 \geq 5$ \checkmark
  \item Aproximación normal es válida
\end{itemize}
\end{frame}

\begin{frame}{IC para una Proporción}
\textbf{Problema}: La fórmula teórica requiere conocer $p$, pero eso es lo que queremos estimar.

\vspace{0.3cm}
\textbf{Solución}: Reemplazar $p$ por $\hat{p}$ en el error estándar.

\vspace{0.3cm}
\begin{block}{IC al $(1-\alpha) \times 100\%$ para $p$}
\[
\hat{p} \pm z_{\alpha/2} \sqrt{\frac{\hat{p}(1-\hat{p})}{n}}
\]
Equivalentemente:
\[
\left( \hat{p} - z_{\alpha/2} \sqrt{\frac{\hat{p}(1-\hat{p})}{n}}, \quad \hat{p} + z_{\alpha/2} \sqrt{\frac{\hat{p}(1-\hat{p})}{n}} \right)
\]
\end{block}

\vspace{0.3cm}
\textbf{Nota}: Usamos $z$ (no $t$) porque la fórmula del error estándar ya usa $\hat{p}$ como estimador de $p$.
\end{frame}

\begin{frame}[fragile]{Ejemplo: Proporción de Estudiantes que Superan B1}
\textbf{Contexto}: En los datos ICFES, definimos éxito como \variable{MOD\_INGLES\_PUNT} $\geq 160$ (aproximadamente B1).

\begin{lstlisting}
# Cargar datos
icfes_ni <- read.csv("datos/icfes_ni_2023.csv")

# Crear variable binaria (1 = supera B1, 0 = no supera)
icfes_ni$supera_B1 <- ifelse(icfes_ni$MOD_INGLES_PUNT >= 160, 1, 0)

# Calcular proporcion muestral
n <- nrow(icfes_ni)
x <- sum(icfes_ni$supera_B1, na.rm = TRUE)
p_hat <- x / n

cat("Tamano de muestra:", n)
cat("\nExitos (supera B1):", x)
cat("\nProporcion muestral:", round(p_hat, 4))

# Verificar condicion de aplicabilidad
cat("\nn*p_hat =", n * p_hat)
cat("\nn*(1-p_hat) =", n * (1 - p_hat))
\end{lstlisting}
\end{frame}

\begin{frame}[fragile]{Calcular IC Manualmente}
\begin{lstlisting}
# IC al 95% usando la formula
alpha <- 0.05
z_critico <- qnorm(1 - alpha/2)

# Error estandar
ee <- sqrt(p_hat * (1 - p_hat) / n)

# Margen de error
E <- z_critico * ee

# Limites del IC
ic_inf <- p_hat - E
ic_sup <- p_hat + E

cat("Error estandar:", round(ee, 4))
cat("\nMargen de error:", round(E, 4))
cat("\nIC al 95%: (", round(ic_inf, 4), ",",
    round(ic_sup, 4), ")")
cat("\nEn porcentaje: (", round(ic_inf*100, 2), "%,",
    round(ic_sup*100, 2), "%)")
\end{lstlisting}
\end{frame}

\begin{frame}[fragile]{Usar \code{prop.test()} en R}
\R\ tiene funciones integradas para IC de proporciones:

\begin{lstlisting}
# Metodo 1: prop.test() (usa correccion de continuidad)
resultado_prop <- prop.test(x, n, conf.level = 0.95)
print(resultado_prop)

# Extraer IC
ic_prop <- resultado_prop$conf.int
cat("\nIC con prop.test():", round(ic_prop, 4))

# Metodo 2: binom.test() (metodo exacto, recomendado para n pequeno)
resultado_binom <- binom.test(x, n, conf.level = 0.95)
ic_binom <- resultado_binom$conf.int
cat("\nIC con binom.test():", round(ic_binom, 4))
\end{lstlisting}

\vspace{0.3cm}
\textbf{Diferencias}:
\begin{itemize}
  \item \code{prop.test()}: usa aproximación normal con corrección de continuidad
  \item \code{binom.test()}: usa distribución binomial exacta (mejor para $n$ pequeño)
\end{itemize}
\end{frame}

\begin{frame}{Interpretación del Resultado}
\textbf{Ejemplo hipotético}: $n = 500$, $x = 360$, $\hat{p} = 0.72$

\vspace{0.3cm}
IC al 95\%: $(0.680, 0.760)$ o $(68.0\%, 76.0\%)$

\vspace{0.3cm}
\begin{exampleblock}{Interpretación}
Estamos 95\% confiados de que la proporción poblacional de estudiantes de NI que superan el nivel B1 en inglés está entre 68\% y 76\%.
\end{exampleblock}

\vspace{0.3cm}
\textbf{Implicación práctica}:
\begin{itemize}
  \item Si la meta institucional es que al menos 75\% supere B1, no podemos afirmar con 95\% de confianza que se cumple (porque 75\% está en el límite superior)
  \item Necesitaríamos más datos o un nivel de confianza menor
\end{itemize}
\end{frame}

\begin{frame}[fragile]{IC por Departamento}
\textbf{Pregunta}: ¿Varía la proporción de estudiantes que superan B1 entre departamentos?

\begin{lstlisting}
library(dplyr)

# Calcular IC de proporcion por departamento
ic_prop_depto <- icfes_ni %>%
  group_by(ESTU_DEPTO_RESIDE) %>%
  summarise(
    n = n(),
    exitos = sum(supera_B1, na.rm = TRUE),
    p_hat = exitos / n,
    ee = sqrt(p_hat * (1 - p_hat) / n),
    z_critico = qnorm(0.975),
    ic_inf = p_hat - z_critico * ee,
    ic_sup = p_hat + z_critico * ee
  ) %>%
  filter(n >= 30)  # Solo deptos con muestra suficiente

# Ordenar por p_hat
ic_prop_depto <- ic_prop_depto %>%
  arrange(desc(p_hat))

head(ic_prop_depto, 10)
\end{lstlisting}
\end{frame}

\begin{frame}[fragile]{Forest Plot de Proporciones}
\begin{lstlisting}
library(ggplot2)

ggplot(ic_prop_depto,
       aes(x = p_hat, y = reorder(ESTU_DEPTO_RESIDE, p_hat))) +
  geom_point(size = 3, color = "darkgreen") +
  geom_errorbarh(aes(xmin = ic_inf, xmax = ic_sup),
                 height = 0.3, color = "darkgreen") +
  geom_vline(xintercept = weighted.mean(ic_prop_depto$p_hat,
                                         ic_prop_depto$n),
             linetype = "dashed", color = "red") +
  scale_x_continuous(labels = scales::percent) +
  labs(x = "Proporcion que supera B1 (IC 95%)",
       y = "Departamento",
       title = "Proporcion de Estudiantes de NI que Superan B1 en Ingles",
       subtitle = "Por departamento (n >= 30)") +
  theme_minimal()
\end{lstlisting}
\end{frame}

\begin{frame}[fragile]{Comparar Dos Proporciones}
\textbf{Pregunta}: ¿Difiere la proporción de estudiantes que superan B1 entre IES públicas y privadas?

\begin{lstlisting}
# Crear tabla de contingencia
tabla <- table(icfes_ni$TIPO_IES, icfes_ni$supera_B1)
print(tabla)

# Prueba de proporciones (adelanto de pruebas de hipotesis)
test_2prop <- prop.test(tabla, conf.level = 0.95)
print(test_2prop)

# Calcular IC separados
ic_publica <- prop.test(tabla[1,2], sum(tabla[1,]), conf.level = 0.95)
ic_privada <- prop.test(tabla[2,2], sum(tabla[2,]), conf.level = 0.95)

cat("IC publica:", round(ic_publica$conf.int, 4))
cat("\nIC privada:", round(ic_privada$conf.int, 4))
\end{lstlisting}

\vspace{0.3cm}
\small
\textbf{Nota}: Si los dos IC NO se solapan, hay evidencia de diferencia significativa.
\end{frame}

% ============================================================================
\section{Tamaño de Muestra para Medias}
% ============================================================================

\begin{frame}{Planificación del Muestreo}
\textbf{Pregunta inversa}: En lugar de calcular el IC dado $n$, ¿cuál debe ser $n$ para obtener un IC con un margen de error deseado?

\vspace{0.3cm}
\textbf{Aplicaciones}:
\begin{itemize}
  \item Encuesta de satisfacción: ¿cuántos estudiantes de NI encuestar para estimar el salario inicial promedio con error $\pm \$300{,}000$ COP?
  \item Estudio de mercado: ¿cuántos usuarios de Rappi necesito para estimar el tiempo promedio de entrega con precisión de $\pm 2$ minutos?
  \item Asignar presupuesto: más $n$ $\Rightarrow$ mayor costo de recolección
\end{itemize}

\vspace{0.3cm}
\textbf{Componentes del problema}:
\begin{itemize}
  \item Margen de error deseado $E$
  \item Nivel de confianza $(1-\alpha)$
  \item Estimación previa de la variabilidad ($\sigma$ o $s$)
\end{itemize}
\end{frame}

\begin{frame}{Fórmula de Tamaño de Muestra para $\mu$}
Recordemos que el margen de error es:
\[
E = z_{\alpha/2} \frac{\sigma}{\sqrt{n}}
\]

Despejando $n$:
\[
n = \left( \frac{z_{\alpha/2} \cdot \sigma}{E} \right)^2
\]

\vspace{0.3cm}
\begin{block}{Tamaño de muestra para estimar $\mu$}
\[
n = \left( \frac{z_{\alpha/2} \cdot \sigma}{E} \right)^2
\]
donde:
\begin{itemize}
  \item $E$ = margen de error deseado
  \item $\sigma$ = desviación estándar poblacional (estimada de un estudio piloto o literatura)
\end{itemize}
\textbf{Redondear siempre hacia arriba} para garantizar la precisión deseada.
\end{block}
\end{frame}

\begin{frame}[fragile]{Ejemplo: Tamaño de Muestra para MOD\_INGLES\_PUNT}
\textbf{Escenario}: Queremos estimar la media de \variable{MOD\_INGLES\_PUNT} en NI con un margen de error de $\pm 2$ puntos al 95\% de confianza. De estudios previos, sabemos que $\sigma \approx 18$.

\begin{lstlisting}
# Parametros
E <- 2           # Margen de error deseado
sigma <- 18      # Desviacion estandar estimada
confianza <- 0.95
alpha <- 1 - confianza

# Valor critico
z_critico <- qnorm(1 - alpha/2)

# Calcular n
n <- (z_critico * sigma / E)^2

cat("z critico:", round(z_critico, 4))
cat("\nn calculado:", n)
cat("\nn redondeado hacia arriba:", ceiling(n))
\end{lstlisting}

\vspace{0.3cm}
\textbf{Resultado}: Necesitamos $n = \lceil 311.17 \rceil = 312$ estudiantes.
\end{frame}

\begin{frame}{Relación entre $n$, $E$, y Nivel de Confianza}
\textbf{Efecto de diferentes parámetros}:

\vspace{0.3cm}
\begin{enumerate}
  \item \textbf{Mayor precisión (menor $E$)} $\Rightarrow$ mayor $n$
  \begin{itemize}
    \item Si queremos $E = 1$ en lugar de $E = 2$, entonces $n$ se multiplica por $2^2 = 4$
    \item $n \propto 1/E^2$
  \end{itemize}

  \item \textbf{Mayor confianza} $\Rightarrow$ mayor $n$
  \begin{itemize}
    \item 95\% de confianza: $z = 1.96$
    \item 99\% de confianza: $z = 2.576$ ($+31\%$ más grande)
    \item $n$ aumenta en $(2.576/1.96)^2 = 1.73$, es decir, 73\% más
  \end{itemize}

  \item \textbf{Mayor variabilidad ($\sigma$)} $\Rightarrow$ mayor $n$
  \begin{itemize}
    \item $n \propto \sigma^2$
  \end{itemize}
\end{enumerate}
\end{frame}

\begin{frame}[fragile]{Sensibilidad del Tamaño de Muestra}
\begin{lstlisting}
# Explorar diferentes combinaciones
library(tidyr)

combinaciones <- expand.grid(
  E = c(1, 2, 3, 5),
  confianza = c(0.90, 0.95, 0.99)
)

combinaciones$z <- qnorm(1 - (1 - combinaciones$confianza)/2)
combinaciones$n <- ceiling((combinaciones$z * sigma / combinaciones$E)^2)

# Tabla pivotada
tabla_n <- combinaciones %>%
  select(E, confianza, n) %>%
  pivot_wider(names_from = confianza, values_from = n,
              names_prefix = "Conf_")

print(tabla_n)
\end{lstlisting}

\vspace{0.3cm}
\small
\textbf{Interpretación}: A medida que $E$ disminuye o la confianza aumenta, $n$ crece rápidamente.
\end{frame}

\begin{frame}[fragile]{¿Qué Hacer si No Conocemos $\sigma$?}
\textbf{Opciones}:

\vspace{0.3cm}
\begin{enumerate}
  \item \textbf{Estudio piloto}: Tomar una muestra pequeña ($n_{\text{piloto}} = 30-50$), calcular $s$ y usar como estimación de $\sigma$.

  \item \textbf{Literatura}: Buscar estudios previos similares y usar su $\sigma$ reportada.

  \item \textbf{Rango}: Si conocemos el rango aproximado de la variable, usar $\sigma \approx \text{rango}/4$ o $\text{rango}/6$.
\end{enumerate}

\begin{lstlisting}
# Ejemplo: estudio piloto
piloto <- sample(icfes_ni$MOD_INGLES_PUNT, 50, replace = FALSE)
sigma_estimado <- sd(piloto, na.rm = TRUE)

cat("Sigma estimado del piloto:", round(sigma_estimado, 2))

# Usar sigma_estimado para calcular n
n_final <- ceiling((qnorm(0.975) * sigma_estimado / 2)^2)
cat("\nTamano de muestra necesario:", n_final)
\end{lstlisting}
\end{frame}

% ============================================================================
\section{Tamaño de Muestra para Proporciones}
% ============================================================================

\begin{frame}{Fórmula de Tamaño de Muestra para $p$}
Recordemos que el margen de error para una proporción es:
\[
E = z_{\alpha/2} \sqrt{\frac{p(1-p)}{n}}
\]

Despejando $n$:
\[
n = \frac{z_{\alpha/2}^2 \cdot p(1-p)}{E^2}
\]

\vspace{0.3cm}
\begin{alertblock}{Problema}
La fórmula requiere conocer $p$, pero eso es lo que queremos estimar.
\end{alertblock}

\vspace{0.3cm}
\textbf{Soluciones}:
\begin{itemize}
  \item Usar una estimación previa $p^*$ (de estudios piloto o literatura)
  \item Usar $p^* = 0.5$ (enfoque conservador)
\end{itemize}
\end{frame}

\begin{frame}{Enfoque Conservador: $p^* = 0.5$}
\textbf{Propiedad}: El producto $p(1-p)$ alcanza su máximo cuando $p = 0.5$.

\begin{center}
\begin{tikzpicture}[scale=0.8]
  \draw[->] (0,0) -- (5.5,0) node[right] {$p$};
  \draw[->] (0,0) -- (0,3) node[above] {$p(1-p)$};

  \draw[domain=0:5, smooth, variable=\x, thick, blue, samples=100]
    plot ({\x},{12*(\x/5)*(1-\x/5)});

  \draw[dashed] (2.5,0) -- (2.5,3);
  \node[below] at (2.5,0) {0.5};
  \node[above] at (2.5,3) {max = 0.25};

  \node[below] at (0,0) {0};
  \node[below] at (5,0) {1};
\end{tikzpicture}
\end{center}

\vspace{0.3cm}
\textbf{Implicación}: Usar $p^* = 0.5$ garantiza que $n$ será suficiente \alert{sin importar el valor verdadero de $p$}.

\vspace{0.3cm}
\begin{block}{Fórmula conservadora}
\[
n = \frac{z_{\alpha/2}^2 \cdot 0.25}{E^2} = \frac{z_{\alpha/2}^2}{4E^2}
\]
\end{block}
\end{frame}

\begin{frame}[fragile]{Ejemplo: Tamaño de Muestra para Proporción}
\textbf{Escenario}: Queremos estimar la proporción de estudiantes de NI que superan B2 en inglés (\variable{MOD\_INGLES\_PUNT} $\geq 180$) con un margen de error de $\pm 3\%$ al 95\% de confianza.

\begin{lstlisting}
# Parametros
E <- 0.03        # Margen de error (3%)
confianza <- 0.95
z_critico <- qnorm(1 - (1 - confianza)/2)

# Metodo 1: Enfoque conservador (p* = 0.5)
n_conservador <- ceiling((z_critico^2 * 0.25) / E^2)
cat("n conservador (p*=0.5):", n_conservador)

# Metodo 2: Estimacion previa (supongamos p* = 0.20)
p_star <- 0.20
n_estimado <- ceiling((z_critico^2 * p_star * (1 - p_star)) / E^2)
cat("\nn con p*=0.20:", n_estimado)

# Diferencia
cat("\nAhorro de muestra:", n_conservador - n_estimado)
\end{lstlisting}
\end{frame}

\begin{frame}{Comparación de Métodos}
\textbf{Resultados del ejemplo}:
\begin{itemize}
  \item Con $p^* = 0.5$ (conservador): $n = 1068$
  \item Con $p^* = 0.20$ (estimado): $n = 683$
  \item Ahorro: $1068 - 683 = 385$ estudiantes (36\% menos)
\end{itemize}

\vspace{0.3cm}
\textbf{Trade-off}:
\begin{itemize}
  \item \textbf{Conservador}: Garantiza precisión sin importar $p$, pero puede ser costoso
  \item \textbf{Estimado}: Más eficiente, pero arriesgado si $p^*$ está muy lejos de $p$
\end{itemize}

\vspace{0.3cm}
\begin{exampleblock}{Recomendación}
Si tienes una estimación razonable de $p$ (de un piloto o literatura), úsala. Si no, usa $p^* = 0.5$ para estar seguro.
\end{exampleblock}
\end{frame}

\begin{frame}[fragile]{Función en R para Tamaño de Muestra}
\begin{lstlisting}
# Funcion personalizada
calcular_n_proporcion <- function(E, confianza = 0.95, p_star = 0.5) {
  z <- qnorm(1 - (1 - confianza)/2)
  n <- (z^2 * p_star * (1 - p_star)) / E^2
  return(ceiling(n))
}

# Ejemplos de uso
cat("n para E=0.03, conf=95%, p*=0.5:",
    calcular_n_proporcion(0.03))

cat("\nn para E=0.05, conf=95%, p*=0.5:",
    calcular_n_proporcion(0.05))

cat("\nn para E=0.03, conf=99%, p*=0.5:",
    calcular_n_proporcion(0.03, confianza = 0.99))

cat("\nn para E=0.03, conf=95%, p*=0.20:",
    calcular_n_proporcion(0.03, p_star = 0.20))
\end{lstlisting}
\end{frame}

% ============================================================================
\section{Muestreo Estratificado y Consideraciones Prácticas}
% ============================================================================

\begin{frame}{Muestreo Aleatorio Simple vs Estratificado}
\begin{columns}[T]
\column{0.48\textwidth}
\textbf{Muestreo Aleatorio Simple (MAS)}:
\begin{itemize}
  \item Cada elemento tiene igual probabilidad de selección
  \item Fácil de implementar
  \item Puede no ser representativo de subgrupos
\end{itemize}

\vspace{0.3cm}
\textbf{Ejemplo}: Seleccionar 500 estudiantes de NI al azar de toda Colombia.

\column{0.48\textwidth}
\textbf{Muestreo Estratificado}:
\begin{itemize}
  \item Dividir la población en \alert{estratos} homogéneos
  \item Muestrear dentro de cada estrato
  \item Garantiza representatividad
\end{itemize}

\vspace{0.3cm}
\textbf{Ejemplo}: Dividir por departamento, luego seleccionar proporcionalmente de cada uno.
\end{columns}

\vspace{0.5cm}
\begin{block}{Estratos comunes en educación}
\begin{itemize}
  \item Departamento o región
  \item Tipo de IES (pública/privada)
  \item Nivel socioeconómico
  \item Género
\end{itemize}
\end{block}
\end{frame}

\begin{frame}{Ventajas del Muestreo Estratificado}
\textbf{1. Mayor precisión}:
\begin{itemize}
  \item Si los estratos son homogéneos internamente, la variabilidad dentro de cada estrato es menor que la variabilidad global
  \item Esto reduce el error estándar de las estimaciones
\end{itemize}

\vspace{0.3cm}
\textbf{2. Representatividad garantizada}:
\begin{itemize}
  \item Asegura que cada subgrupo importante esté representado
  \item Evita muestras "raras" (por ejemplo, todos de Bogotá)
\end{itemize}

\vspace{0.3cm}
\textbf{3. Estimaciones por subgrupo}:
\begin{itemize}
  \item Permite calcular IC separados para cada estrato
  \item Útil para comparaciones regionales
\end{itemize}

\vspace{0.3cm}
\textbf{4. Eficiencia administrativa}:
\begin{itemize}
  \item Facilita la logística de recolección de datos
\end{itemize}
\end{frame}

\begin{frame}{Asignación Proporcional vs Óptima}
\textbf{Asignación proporcional}:
\[
n_i = n \cdot \frac{N_i}{N}
\]
donde $N_i$ = tamaño del estrato $i$, $N$ = tamaño poblacional total.

\vspace{0.3cm}
\textbf{Ejemplo}: Si el 15\% de los estudiantes de NI están en Antioquia, asignamos el 15\% de la muestra a Antioquia.

\vspace{0.3cm}
\textbf{Asignación óptima (Neyman)}:
\[
n_i = n \cdot \frac{N_i \sigma_i}{\sum_{j=1}^k N_j \sigma_j}
\]
Asigna más muestra a estratos con \alert{mayor variabilidad} $\sigma_i$.

\vspace{0.3cm}
\begin{alertblock}{Limitación}
La asignación óptima requiere conocer $\sigma_i$ de antemano (estudios piloto o literatura).
\end{alertblock}
\end{frame}

\begin{frame}[fragile]{Ejemplo: Asignación Proporcional por Departamento}
\begin{lstlisting}
# Supongamos que queremos n = 1000 estudiantes de NI

# Paso 1: Conocer la distribucion poblacional
poblacion_depto <- icfes_ni %>%
  count(ESTU_DEPTO_RESIDE, name = "N_i") %>%
  mutate(peso = N_i / sum(N_i))

# Paso 2: Asignar muestra proporcionalmente
n_total <- 1000
poblacion_depto$n_i <- ceiling(n_total * poblacion_depto$peso)

# Ajuste fino: asegurar que sum(n_i) = n_total
exceso <- sum(poblacion_depto$n_i) - n_total
if (exceso > 0) {
  # Reducir de los estratos mas grandes
  poblacion_depto$n_i[order(-poblacion_depto$N_i)[1:exceso]] <-
    poblacion_depto$n_i[order(-poblacion_depto$N_i)[1:exceso]] - 1
}

print(poblacion_depto)
\end{lstlisting}
\end{frame}

\begin{frame}{Consideraciones Prácticas en el Diseño Muestral}
\textbf{1. Costo vs Precisión}:
\begin{itemize}
  \item Mayor $n$ $\Rightarrow$ mayor precisión, pero también mayor costo
  \item Equilibrio: ¿cuánta precisión necesitamos realmente?
\end{itemize}

\vspace{0.3cm}
\textbf{2. Tasa de respuesta}:
\begin{itemize}
  \item No todos los seleccionados responderán
  \item Ajustar: si esperamos 70\% de respuesta y queremos $n = 500$, seleccionar $500/0.7 \approx 715$
\end{itemize}

\vspace{0.3cm}
\textbf{3. Datos faltantes}:
\begin{itemize}
  \item Aumentar $n$ para compensar valores perdidos
  \item Métodos de imputación
\end{itemize}

\vspace{0.3cm}
\textbf{4. Poder estadístico}:
\begin{itemize}
  \item Para pruebas de hipótesis (próximas sesiones), $n$ debe ser suficiente para \alert{detectar} diferencias relevantes
\end{itemize}
\end{frame}

% ============================================================================
\section{Resumen de toda la Unidad 2}
% ============================================================================

\begin{frame}{Tabla Maestra de Intervalos de Confianza}
\begin{table}
\centering
\footnotesize
\begin{tabular}{p{2cm}p{3.5cm}p{4.5cm}}
\toprule
\textbf{Parámetro} & \textbf{Condiciones} & \textbf{IC al $(1-\alpha) \times 100\%$} \\
\midrule
$\mu$ & $\sigma$ conocida, $n$ grande o población normal & $\bar{x} \pm z_{\alpha/2} \frac{\sigma}{\sqrt{n}}$ \\[0.3cm]
$\mu$ & $\sigma$ desconocida, población normal o $n \geq 30$ & $\bar{x} \pm t_{\alpha/2,n-1} \frac{s}{\sqrt{n}}$ \\[0.3cm]
$\sigma^2$ & Población normal & $\left(\frac{(n-1)s^2}{\chi^2_{\alpha/2}}, \frac{(n-1)s^2}{\chi^2_{1-\alpha/2}}\right)$ \\[0.3cm]
$p$ & $n\hat{p} \geq 5$ y $n(1-\hat{p}) \geq 5$ & $\hat{p} \pm z_{\alpha/2} \sqrt{\frac{\hat{p}(1-\hat{p})}{n}}$ \\
\bottomrule
\end{tabular}
\end{table}
\end{frame}

\begin{frame}{Tabla de Tamaño de Muestra}
\begin{table}
\centering
\small
\begin{tabular}{llp{5cm}}
\toprule
\textbf{Parámetro} & \textbf{Fórmula} & \textbf{Notas} \\
\midrule
$\mu$ & $n = \left(\frac{z_{\alpha/2} \cdot \sigma}{E}\right)^2$ & Requiere estimación previa de $\sigma$ \\[0.3cm]
$p$ & $n = \frac{z_{\alpha/2}^2 \cdot p^*(1-p^*)}{E^2}$ & Usar $p^* = 0.5$ si no hay estimación previa \\
\bottomrule
\end{tabular}
\end{table}

\vspace{0.3cm}
\textbf{Recordar}: Siempre redondear $n$ hacia arriba.
\end{frame}

\begin{frame}{Conexión con Pruebas de Hipótesis}
\textbf{Adelanto de la próxima unidad}:

\vspace{0.3cm}
\begin{block}{Relación IC - Pruebas de hipótesis}
Si un IC al $(1-\alpha) \times 100\%$ NO contiene un valor $\theta_0$, entonces rechazamos $H_0: \theta = \theta_0$ al nivel de significancia $\alpha$.
\end{block}

\vspace{0.3cm}
\textbf{Ejemplo}:
\begin{itemize}
  \item IC al 95\% para $\mu$: $(162, 168)$
  \item $H_0: \mu = 160$ vs $H_1: \mu \neq 160$
  \item Como 160 NO está en $(162, 168)$, rechazamos $H_0$ al nivel $\alpha = 0.05$
\end{itemize}

\vspace{0.3cm}
\textbf{Ventaja de los IC}: Proporcionan más información que solo "rechazar" o "no rechazar" $H_0$. Dan un rango de valores plausibles.
\end{frame}

\begin{frame}{Resumen Conceptual de la Unidad 2}
\textbf{Lógica general de la inferencia}:

\vspace{0.3cm}
\begin{enumerate}
  \item \textbf{Muestra}: Observamos datos de una muestra aleatoria
  \item \textbf{Estadístico}: Calculamos un estimador puntual ($\bar{x}$, $\hat{p}$, $s^2$)
  \item \textbf{Distribución muestral}: Sabemos cómo se comporta el estadístico (TLC, distribuciones $t$, $\chi^2$)
  \item \textbf{IC}: Construimos un intervalo que probablemente contiene el parámetro
  \item \textbf{Interpretación}: Comunicamos incertidumbre de forma cuantitativa
\end{enumerate}

\vspace{0.3cm}
\textbf{Próximos pasos}:
\begin{itemize}
  \item Unidad 3: Pruebas de hipótesis (¿es $\mu$ significativamente diferente de un valor?)
  \item Unidad 4: Regresión lineal (¿cómo se relacionan dos variables?)
\end{itemize}
\end{frame}

% ============================================================================
\section{Conclusiones}
% ============================================================================

\begin{frame}{Resumen de la Sesión}
\textbf{Hoy aprendimos}:
\begin{enumerate}
  \item \textbf{IC para proporciones}:
  \begin{itemize}
    \item $\hat{p} \pm z_{\alpha/2} \sqrt{\frac{\hat{p}(1-\hat{p})}{n}}$
    \item Condición: $n\hat{p} \geq 5$ y $n(1-\hat{p}) \geq 5$
    \item Funciones en R: \code{prop.test()}, \code{binom.test()}
  \end{itemize}

  \item \textbf{Tamaño de muestra para medias}:
  \begin{itemize}
    \item $n = \left(\frac{z_{\alpha/2} \cdot \sigma}{E}\right)^2$
  \end{itemize}

  \item \textbf{Tamaño de muestra para proporciones}:
  \begin{itemize}
    \item $n = \frac{z_{\alpha/2}^2 \cdot p^*(1-p^*)}{E^2}$
    \item Enfoque conservador: $p^* = 0.5$
  \end{itemize}

  \item \textbf{Muestreo estratificado}: Mayor precisión y representatividad
\end{enumerate}

\vspace{0.3cm}
\textbf{Próxima sesión}: Fundamentos de pruebas de hipótesis ($H_0$, $H_1$, p-valor, errores tipo I y II).
\end{frame}

\begin{frame}{Ejercicios para la Próxima Clase}
\textbf{Con los datos de ICFES}:
\begin{enumerate}
  \item Calcular IC al 95\% para la proporción de estudiantes de NI que superan B2 en inglés (\variable{MOD\_INGLES\_PUNT} $\geq 180$)
  \item Crear forest plot de proporciones por región (Caribe, Andina, Pacífico, etc.)
  \item Determinar el tamaño de muestra necesario para estimar la media de \variable{MOD\_RAZONA\_CUANTITAT} con error $\pm 3$ puntos al 99\% de confianza
  \item Comparar IC de proporción de aprobación (puntaje global $\geq 300$) entre IES públicas y privadas
  \item \textbf{Desafío}: Diseñar un muestreo estratificado por departamento con asignación proporcional para $n = 800$
\end{enumerate}

\vspace{0.3cm}
\textbf{Lectura}: Anderson et al., Capítulos 8-9 (Estimación e Intervalos de Confianza)
\end{frame}

\begin{frame}
\begin{center}
\Large
¡Gracias por su atención!

\vspace{1cm}
\normalsize
Preguntas y discusión
\end{center}
\end{frame}

\end{document}
